\chapter{总览}

\begin{quote}
机遇偏爱有准备的头脑。

--- 路易斯$\cdot$巴斯德
\end{quote}

\section{引言}

本章引入主动推断试图回答的核心问题：生命有机体如何在与环境进行适应性交换的同时持续存在？我们将讨论为何要从规范性视角来处理这一问题。这样的视角从第一性原理出发，再逐步展开其在认知和生物学上的含义。此外，本章还将简要介绍全书结构，包括它被划分为两个部分：第一部分旨在帮助读者理解主动推断，第二部分旨在帮助读者将其用于自己的研究。

\section{生命有机体如何持续存在并进行适应性行动？}

生命有机体持续不断地与其环境（包括其他有机体）发生相互作用。它们发出会改变环境的动作，并从环境中接收感觉观察，这一点如图 1.1 的示意所示。

生命有机体只有对行动--感知回路施加适应性控制，才能维持其身体完整性。这意味着，它们要通过行动来引出那些要么对应于理想结果或目标、要么有助于理解世界的感觉观察。前者例如对简单有机体而言，获得营养和庇护时所伴随的感觉；对更复杂的有机体而言，则可能是朋友和工作。后者则例如为有机体提供关于其周围环境的信息。

与环境进入适应性的行动--感知回路，会给生命有机体带来严峻挑战。这在很大程度上源于该循环的递归性质：每一次由前一个动作引出的观察，都会改变我们如何决定下一个动作，从而再去引出下一个观察。控制与适应的可能性极其丰富，但真正有用的却很少。尽管如此，在进化过程中，生命有机体还是发展出了适应性策略，以应对生存的根本挑战。这些策略在认知复杂度上彼此不同：较简单的有机体采用更简单、更刚性的解决方案（例如细菌沿着营养梯度移动），而更高级的有机体则采用认知要求更高、也更灵活的解决方案（例如人类为实现远端目标而进行规划）。这些策略在被选择和发挥作用的时间尺度上也不相同：有的体现为对环境威胁的简单反应，或是在进化时间尺度上形成的形态适应；有的体现为在文化学习或个体发育学习中建立起来的行为模式；还有的则需要与行动和感知处于相近时间尺度上的认知过程，例如注意和记忆。

\section{主动推断：从第一性原理推导行为}

这种多样性对生物学而言是一种福音，但对大脑与心智的形式理论而言却是一项挑战。大体上，我们可以采取两种看法。

一种看法认为，不同的生物适应、神经过程（例如突触交换与脑网络）以及认知机制（例如感知、注意、社会互动）都具有很强的特殊性，因此需要各自专门的解释。这样一来，哲学、心理学、神经科学、动物行为学、生物学、人工智能和机器人学等领域中的理论就会不断繁殖，而几乎看不到统一它们的希望。

另一种看法认为，尽管生命有机体中的行为、认知和适应具有多样化的表现形式，但它们的核心方面仍然可以从第一性原理出发获得连贯解释。

这两种可能性对应着两种不同的研究纲领，在某种程度上也对应着两种不同的科学态度：``整洁派''与``杂糅派''（这两个术语出自 Roger Shank）。整洁派总是试图在大脑与心智现象表面的异质性背后寻求统一。这通常对应于构建自上而下的规范性模型\footnote{这里的``规范性''是指：模型从系统应当满足的原则或约束出发，而不是仅仅从经验数据中做归纳。}，从第一性原理起步，并尝试尽可能多地推导出关于大脑与心智的结论。相反，杂糅派则拥抱异质性，聚焦于那些要求专门解释的细节。这通常对应于构建自下而上的模型：从数据出发，采用任何有效的方法去解释复杂现象，甚至允许不同现象对应不同解释。

那么，是否可能像整洁派所设想的那样，从第一性原理解释异质的生物学与认知现象？是否可能建立一个统一的框架来理解大脑与心智？

本书对这些问题给出肯定回答，并将主动推断推进为一种理解大脑与心智的规范性路径。我们对主动推断的处理从第一性原理出发，并逐步展开它们在认知与生物学上的含义。

\section{本书结构}

本书由两个部分组成。它们分别面向两类读者：一类希望理解主动推断（第一部分），另一类希望将其用于自己的研究（第二部分）。本书第一部分将从概念和形式两个层面介绍主动推断，并把它放入当前认知理论的背景中加以定位。其目标是对主动推断给出一个全面、形式化且自洽的入门介绍：包括它的主要构件，以及它对于大脑与认知研究的含义。

本书第二部分将展示若干使用主动推断来解释认知现象的计算模型实例，例如感知、注意、记忆和规划。其目标是帮助读者既能理解现有的主动推断计算模型，也能设计新的模型。简言之，本书分为理论（第一部分）与实践（第二部分）。

\subsection{第一部分：理论中的主动推断}

主动推断是一个规范性框架，用于刻画生命有机体中的贝叶斯最优\footnote{这里的``贝叶斯最优''是指：在给定概率模型与目标偏好的前提下，以概率推断和决策原则达到最优。}行为与认知。它的规范性特征体现在这样一个观点中：生命有机体中行为与认知的所有方面，都服从同一个唯一的命令，即最小化其感觉观察的惊讶（surprise）。这里的惊讶必须从技术意义上理解：它衡量的是，一个主体当前的感觉观察与其偏好的感觉观察之间有多大差异。所谓偏好的感觉观察，就是那些能够维持其完整性的观察（例如，对一条鱼来说，处于水中）。重要的是，最小化惊讶并不是靠被动地观察环境就能做到的；相反，主体必须对其行动--感知回路进行适应性控制，以便引出理想的感觉观察。这正是主动推断中``主动''二字的含义。

由于一些稍后会显现出来的技术原因，最小化惊讶是一个很困难的问题。主动推断对此提供了一种解决方案。它假定，即使生命有机体不能直接最小化惊讶，它们也可以最小化一个代理量，即（变分）自由能。这个量可以通过神经计算，在回应感觉观察以及预期感觉观察时被最小化。对自由能最小化的强调，也揭示了主动推断与激发它的那个（首要）原理之间的关系：自由能原理（Friston 2009）。

以自由能最小化作为出发点来解释生物现象，乍看之下似乎非常抽象。然而，我们可以从中推导出一系列形式和经验上的含义，并回答认知理论与神经理论中的一系列核心问题。这些问题包括：参与自由能最小化的变量如何可能被编码到神经元群体中；自由能最小化计算如何映射到具体认知过程，例如感知、动作选择与学习；以及当主动推断主体最小化其自由能时，会涌现出何种行为。

从上面的议题列表可以看出，本书主要关注的是生命有机体层面的主动推断与自由能最小化，不论这些有机体较为简单（例如细菌）还是较为复杂（例如人类），以及与之相关的行为、认知、社会和神经过程。做出这一澄清，是为了将我们对主动推断的讨论置于更一般的自由能原理（FEP）之中加以定位。后者讨论的是远远超出神经信息处理范围的、更广泛生物现象与时间尺度上的自由能最小化，从进化、细胞到文化层面皆在其内（Friston, Levin et al. 2015; Isomura and Friston 2018; Palacios, Razi et al. 2020; Veissière et al. 2020），而这些并不属于本书的讨论范围。

我们可以沿着两条道路中的任意一条来引出主动推断：一条是高路（high road），另一条是低路（low road）；见图 1.2。这两条道路为主动推断提供了两个彼此不同、但又高度互补的视角。

高路通向主动推断，是从``生命有机体如何在世界中持续存在并适应性行动''这一问题出发，并将主动推断作为这些问题的规范性解答来加以论证。高路视角有助于理解主动推断的规范性本质：生命有机体为了面对其根本性的存在挑战，必须做什么（最小化其自由能），以及为什么必须这样做（以间接方式最小化其感觉观察的惊讶）。

低路通向主动推断，则从``贝叶斯大脑''这一观念出发。该观念将大脑视为一个推断引擎，试图优化关于其感觉输入原因的概率表征。随后，它把主动推断论证为一种针对该推断问题的特定变分近似；否则，这个推断问题往往是不可处理的，而这种近似还具有一定的生物学合理性。低路视角有助于说明主动推断主体如何最小化其自由能，因此它不仅把主动推断展示为一条原则，也把它展示为一种认知功能及其神经基础的机制性解释（即过程理论）。

在第 2 章中，我们将展开主动推断的低路视角。我们从那些把感知视为统计推断（贝叶斯推断）问题的奠基性理论（Helmholtz 1866）出发，再看到它们在贝叶斯大脑假说（Doya 2007）中的现代形态。我们将看到，为了执行这种（感知）推断，生命有机体必须拥有或体现一个概率生成模型，用以说明其感觉观察是如何被生成的。这个模型编码了关于可观察变量（感觉观察）和不可观察（隐藏）变量的信念（概率分布）。我们还会把这一推断视角从感知扩展到动作选择、规划和学习等问题。

在第 3 章中，我们将展示与之互补的主动推断高路视角。该章将引入自由能原理以及生物有机体最小化惊讶这一命令。进一步地，我们会分析这一原理如何涵盖自组织动力学，以及如何维持一个统计边界或马尔可夫毯（Markov blanket），从而保持与环境的分离。这对于维持生物体的完整性至关重要，也是其自创生（autopoiesis）的核心。

在第 4 章中，我们将更为形式化地展开主动推断。该章以第 2 章对贝叶斯大脑的讨论为线索，阐明第 3 章中``自我证成''（self-evidencing）动力学与变分推断之间的数学关系。此外，该章还会介绍两类用于表述主动推断问题的生成模型：一类是用于决策与规划的部分可观察马尔可夫决策过程，另一类是与感觉受体和肌肉相接口的连续时间动力学模型。最后，我们将看到，对于这两类模型，自由能最小化如何体现为动态的信念更新。

在第 5 章中，我们将从形式处理转向主动推断的生物学含义。以``大脑中一切发生变化的东西都必须最小化自由能''（Friston 2009）这一前提出发，我们将讨论参与自由能最小化的具体量（例如预测、预测误差与精度信号）如何体现在神经动力学之中。这有助于把主动推断的抽象计算原理映射到可由生理基质执行的具体神经计算上。这一点对于在该框架下形成假设至关重要，并能确保这些假设可以通过测量数据来回答。换句话说，第 5 章将阐述与主动推断相关联的过程理论。

在本书第一部分中，我们还会持续讨论主动推断的若干典型特征。这些特征凸显了它与其他试图解释生物调节与认知的框架之间的差异，下面先作简要预告。

在主动推断之下，感知与行动是完成同一命令的两种互补方式：最小化自由能。感知通过（贝叶斯）信念更新，也就是改变你的想法，使你的信念与感觉观察相容，从而最小化自由能（以及惊讶）。相反，行动则是通过改变世界，使其与您的信念和目标更加相容，从而最小化自由能（以及惊讶）。这种对认知功能的统一处理，是主动推断与那些将行动和感知彼此割裂开来对待的路径之间的根本差别。学习是最小化自由能的另一种方式，但它在本质上并不不同于感知，只是作用于更慢的时间尺度上。感知与行动之间的互补性将在第 2 章展开。

除了在当下驱动动作选择以改变当前可获得的感觉数据之外，主动推断框架还容纳了规划，也就是在未来选择最优行动路线（或策略）。这里的``最优''是相对于期望自由能来衡量的，它不同于上文在行动与感知背景下讨论的变分自由能。事实上，计算变分自由能依赖于当前和过去的观察，而计算期望自由能还要求对未来观察作出预测（因此称为``期望''）。有趣的是，一个策略的期望自由能由两个部分构成：第一部分量化该策略预计将在多大程度上消除不确定性（探索），第二部分量化预测结果与主体目标的一致程度（利用）。与其他框架不同，主动推断中的策略选择会自动平衡探索与利用。变分自由能与期望自由能之间的关系将在第 2 章展开。

在主动推断之下，所有认知操作都被概念化为对生成模型的推断，这与``大脑执行概率计算''这一观点保持一致，也就是贝叶斯大脑假说。然而，对某一种近似形式的贝叶斯推断的诉求，即一种由第一性原理所驱动的变分方案，则为该过程理论增添了具体性。此外，主动推断把推断路径扩展到了认知中那些很少被纳入考虑的领域，并对生物大脑可能实现的模型类型和推断过程增加了更具体的限定。在某些假设下，主动推断中所用生成模型涌现出来的动力学，与计算神经科学中的一些常见模型高度对应，例如预测编码（Rao and Ballard 1999）和 Helmholtz machine（Dayan et al. 1995）。这一变分方案的细节将在第 4 章展开。

在主动推断之下，感知和学习二者都是主动过程，原因有二。第一，大脑本质上是一台预测机器，它不断预测即将到来的刺激，而不是被动等待刺激到来。这一点很重要，因为感知与学习过程总是被先前的预测所情境化，例如，预期刺激与非预期刺激会以不同方式影响感知与学习。第二，从事主动推断的生物体会主动寻求那些能够消解其不确定性的显著性感觉观察，例如转动其感受器，或者选择信息量更大的学习情节。感知与学习的这种主动特征，与大多数把它们视为大体被动过程的现代理论形成对照；这一点将在第 2 章展开。

行动在本质上是目标导向且带有意图性的。它从一个期望结果或目标出发（类似于控制论中的设定点 set-point 概念），而这一目标被编码为一个先验预测。规划则通过推断出一条能够满足这一预测的动作序列来进行，或者等价地说，通过减少先验预测与当前状态之间的任何预测误差来进行。主动推断中行动的目标导向特征，与早期控制论表述是一致的，但不同于大多数用刺激--反应映射或状态--动作策略来解释行为的现代理论。刺激--反应式或习惯性行为，因此就成为主动推断中更广泛策略家族的一个特例。主动推断的目标导向性质将在第 2 章和第 3 章展开。

主动推断中的各种构件，在大脑中都具有可信的生物学对应物。这意味着，一旦人们为手头问题定义了一个具体生成模型，就能够从主动推断作为规范性理论，转向主动推断作为过程理论，并据此提出具体的经验预测。例如，知觉推断与学习分别对应于突触活动的变化和突触效能的变化；预测的精度（在预测编码中）对应于预测误差单元的突触增益；策略的精度则对应于多巴胺能活动。主动推断的一些生物学后果将在第 5 章展开。

\subsection{第二部分：实践中的主动推断}

如果说本书第一部分为读者提供了理解主动推断所需的概念工具和形式工具，那么第二部分则聚焦于实践问题。具体而言，我们希望为读者提供工具，使他们既能理解主动推断对于认知功能（以及功能障碍）的现有模型，也能设计新的模型。为此，我们将讨论若干使用主动推断的具体模型例子。重要的是，主动推断模型可以沿着不同维度变化，例如采用离散时间还是连续时间表述，采用平坦推断还是层级推断。第二部分的结构如下。

在第 6 章中，我们将介绍一套构建主动推断模型的配方。这套配方涵盖设计有效模型所需的关键步骤，包括识别研究对象系统、选择最合适的生成模型形式（例如用于刻画离散时间或连续时间现象），以及确定应当纳入模型的具体变量。因此，这一章将介绍支撑后续各章模型的设计原则。

在第 7 章中，我们将讨论处理离散时间问题的主动推断模型，例如将问题表述为隐马尔可夫模型（HMM）或部分可观察马尔可夫决策过程（POMDP）。我们的例子包括一个知觉加工模型，以及一个离散觅食选择模型，也就是在决策点处向左还是向右转，以获得奖励。我们还将引入信息寻求、学习与新奇寻求等主题，它们都可以在离散时间的主动推断框架下得到处理。

在第 8 章中，我们将讨论处理连续时间问题的主动推断模型，它们使用随机微分方程来表述。这些模型包括知觉模型（如预测编码）、运动控制模型以及序列动力学模型。有趣的是，正是在连续时间表述中，主动推断的一些最具辨识度的预测会显现出来，例如运动生成源于预测的实现，以及注意现象可以理解为精度控制。我们还将介绍同时包含离散时间变量与连续时间变量的混合主动推断模型。这些模型能够同时评估离散选项之间的选择（例如眼跳的目标）以及由该选择所引出的连续运动（例如眼动运动）。

在第 9 章中，我们将说明如何使用主动推断模型来分析行为实验数据。我们将讨论进行基于模型的数据分析所必需的具体步骤，从数据收集，到模型构建，再到模型反演，以支持对单个参与者层面或群体层面数据的分析。

在第 10 章中，我们将讨论主动推断与心理学、神经科学、人工智能以及哲学中其他理论之间的关系。我们还会强调主动推断中最重要、也最能将其与其他理论区分开来的方面。

在附录中，我们将简要讨论理解本书技术性较强部分所需的数学背景，包括泰勒级数近似、变分拉普拉斯、变分法等概念。此外，我们还会以简洁形式给出主动推断中最重要的方程，供读者查阅。

总而言之，本书第二部分将展示如何使用主动推断来构建关于生物学与认知现象的广泛模型，以及如何设计新的模型的方法论。除具体模型本身的兴趣之外，我们更希望我们的处理方式能澄清：使用一个统一的、规范性的框架，以连贯视角来处理生物学与认知现象，究竟有何价值。归根结底，这正是规范性框架真正的吸引力：它提供统一视角和指导性原则，以调和那些表面上彼此断裂的现象。在这里，这些现象包括感知、决策、注意、学习和运动控制，而在任何一本心理学或神经科学手册中，它们通常都会被分置于不同章节。

第二部分所强调的模型，被选择出来是为了尽可能简单地说明一些特定要点。虽然我们覆盖了若干模型与领域，从离散时间决策到连续时间知觉和运动控制，但显然忽略了许多同样有趣的其他模型。现有文献中还存在许多其他主动推断模型，覆盖的领域极其广泛，包括生物自组织与生命起源（Friston 2013）、形态发生（Friston, Levin et al. 2015）、认知机器人学（Pio-Lopez et al. 2016; Sancaktar et al. 2020）、社会动力学与生态位建构（Bruineberg, Rietveld et al. 2018）、突触网络动力学（Palacios, Isomura et al. 2019）、生物网络中的学习（Friston and Herreros 2016），以及创伤后应激障碍（Linson et al. 2020）和惊恐障碍（Maisto, Barca et al. 2021）等精神病理状况。这些模型在许多维度上各不相同：有些与生物学联系更直接，有些则较弱；有些是单主体模型，有些是多主体模型；有些针对适应性推断，有些则针对失适应性推断（例如患者群体中的推断），等等。

这套不断增长的文献，体现了主动推断日益上升的受欢迎程度，也说明它能够被用于极其广泛的领域。本书的目标，是让读者具备理解并在自身研究中使用主动推断的能力，并且有可能进一步探索它那些尚未被预见的潜力。

\section{总结}

本章从规范性视角出发，简要介绍了主动推断这一解释生物学问题的路径，并预告了这一视角的一些含义，而这些含义将在后续章节中展开。此外，本章还强调了本书分为两个部分：分别帮助读者理解主动推断，以及将其用于自己的研究。在接下来的几章中，我们将发展这里勾勒出的低路与高路视角，然后进一步深入到生成模型的结构及其所导出的消息传递。它们共同构成原则层面的主动推断，并为实践层面的主动推断奠定基础。我们希望这些章节能够说服读者：主动推断不仅为理解行为提供了一个统一原则，也为研究自主系统中的行动与感知提供了一条可处理的路径。
